\section{Type Erasure}

% TODO: replace with Jackson's slides

\begin{frame}[fragile]{List vs Deque Motivation}
  \begin{itemize}
    \item \code{std::list} and \code{std::deque} have the methods \code{push\_back}, \code{pop\_front}, \code{front}.
    \item Thus both can be used as a queue.
  \end{itemize}

  In an ideal world, we can imagine \code{std::list} and \code{std::deque}
  inherit from some shared base class \code{base}, with virtual methods \code{push\_back}, allowing for
  \begin{itemize}
    \item \code{base*} pointer that could point to either \code{std::list} or \code{std::deque}.
    \item dynamic dispatch to correct function when calling \code{push\_back}.
  \end{itemize}

  \begin{minted}{cpp}
    /* 2.02 */
    base *q = likes_deques_more ? new std::deque<int>{} : new std::list<int>{};
    // perform some operations on this queue
    q->push_back(67);
    auto x = q->front();
    q->pop_front();
  \end{minted}
\end{frame}

\begin{frame}{Ideal World}
  \includesvg[inkscapelatex=false, width=\textwidth]{assets/inheritance_list_deque.drawio.svg}
\end{frame}

\begin{frame}[fragile]{Wrappers}
  Reality is often disappointing. In reality, \code{std::list} and \code{std::deque} 
  do not inherit from anything.
  There is no way to obtain a base class pointer that can point to both a \code{std::list} and \code{std::deque}.
  The solution is to wrap them with classes that DO shared a base class.

  \begin{minted}{cpp}
    /* 2.03 */
    class base {
    public:
        virtual void push_back(int val) = 0;
        virtual void pop_front() = 0;
        virtual int front() = 0;
        virtual ~base();
    };
  \end{minted}
\end{frame}

\begin{frame}[fragile]{Wrappers}
  \begin{minted}{cpp}
    struct list_wrapper : public base {
        void push_back(int val) override { elems_.push_back(val); }
        void pop_front() override { elems_.pop_front(); }
        int front() override { return elems_.front(); }
        std::list<int> elems_;
    };

    struct deque_wrapper : public base {
        void push_back(int val) override { elems_.push_back(val); }
        void pop_front() override { elems_.pop_front(); }
        int front() override { return elems_.front(); }
        std::deque<int> elems_;
    };
  \end{minted}
\end{frame}

\begin{frame}{Wrappers}
  \includesvg[inkscapelatex=false, width=\textwidth]{assets/type_erasure_manual.drawio.svg}
\end{frame}

\begin{frame}[fragile]{Templated Wrappers}
  Note that we had to write a wrapper for every type.
  We can automate this by using templating.
  \begin{minted}{cpp}
    /* 2.04 */
    template <typename Q>
    struct wrapper : public base {
        void push_back(int val) override { elems_.push_back(val); }
        void pop_front() override { elems_.pop_front(); }
        int front() override { return elems_.front(); }
        Q elems_;
    };
  \end{minted}
\end{frame}

\begin{frame}[fragile]{Wrappers}
  \begin{minted}{cpp}
    // --- Application Code ---

    int main() {
        std::cout << "Do you like deques more than lists?" << std::endl;
        bool likes_deques_more;
        std::cin >> likes_deques_more;

        // choose the underlying container based on a run-time condition
        base *q = likes_deques_more
            ? (base *) new wrapper<std::deque<int>>{}
            : (base *) new wrapper<std::list<int>>{};

        // perform some operations on this queue
        q->push_back(67);
        auto x = q->front();
        q->pop_front();

        delete q;
    }
  \end{minted}
\end{frame}

\begin{frame}[fragile]{Allocators}
  \begin{sqt:definition}
    An \textbf{allocator} is an object that a container uses to acquire and release memory.
    STL containers accept an allocator as a template parameter.
  \end{sqt:definition}

  \begin{minted}{cpp}
    // std::allocator<int> is the default — calls new/delete
    std::vector<int> v1;

    // equivalent to:
    std::vector<int, std::allocator<int>> v2;
  \end{minted}

  The allocator decides \emph{where} memory comes from: the heap, a pre-allocated buffer, a pool, etc.

  TODO: diagram with a grid showing allocated and free spots.
\end{frame}

\begin{frame}[fragile]{Memory Resources}
  \begin{sqt:definition}
    A \textbf{memory resource} is an object that implements an allocation strategy.
    It provides \code{allocate()} and \code{deallocate()} methods.
  \end{sqt:definition}

  Different resources use different strategies:
  \begin{itemize}
    \item \code{malloc\_resource} --- calls \code{std::aligned\_alloc} / \code{std::free}.
    \item \code{free\_list\_resource} --- manages a pre-allocated buffer with a free list.
    \item Pool allocators, arena allocators, stack allocators, \ldots
  \end{itemize}

  Resources have the same interface but \textbf{do not inherit from a common base class}.
\end{frame}

\begin{frame}[fragile]{The Problem}
  Since the allocator is a template parameter, changing the allocator changes the type:
  \begin{minted}{cpp}
    std::vector<int, MallocAllocator<int>>   v1;  // one type
    std::vector<int, FreeListAllocator<int>> v2;  // different type!
  \end{minted}

  \begin{itemize}
    \item You cannot store \code{v1} and \code{v2} in the same container.
    \item You cannot pass them to the same function.
    \item \textcolor{sqt-question}{Sound familiar?} Same problem as \code{std::list} vs \code{std::deque}.
  \end{itemize}

  \textbf{Solution:} apply the same type erasure pattern to memory resources,
  so a single allocator type can use \emph{any} resource at runtime.

  TODO: diagram showing allocator pointing to a resource
\end{frame}

\begin{frame}[fragile]{Memory Resource Interface}
  TODO: this is contained by the allocator. The allocator invokes virtual methods.

  Define an abstract base class for all memory resources:
  \begin{minted}{cpp}
    /* 2.05 */
    class memory_resource {
    public:
        virtual ~memory_resource() = default;
        virtual void* allocate(size_t bytes, size_t align) = 0;
        virtual void  deallocate(void* p, size_t bytes, size_t align) = 0;
        virtual bool  is_equal(const memory_resource& other) const noexcept = 0;
    };
  \end{minted}
\end{frame}

\begin{frame}[fragile]{Memory Resource Wrapper}
  Concrete resources (e.g. \code{malloc\_resource}) do not inherit from \code{memory\_resource}.
  We wrap them with the same templated wrapper pattern:
  \begin{minted}{cpp}
    /* 2.06 */
    template <typename Resource>
    class memory_resource_wrapper : public memory_resource {
    public:
        template <typename... Args>
        explicit memory_resource_wrapper(Args&&... args)
            : resource_(std::forward<Args>(args)...) {}

    private:
        void* allocate(size_t bytes, size_t align) override {
            return resource_.allocate(bytes, align);
        }
        void deallocate(void* p, size_t bytes, size_t align) override {
            resource_.deallocate(p, bytes, align);
        }
        Resource resource_;
    };
  \end{minted}
\end{frame}

\begin{frame}[fragile]{Concrete Resources}
  Concrete resources are plain classes --- no inheritance required.
  \begin{minted}{cpp}
    /* 2.05 */
    class malloc_resource {
    public:
        void* allocate(size_t bytes, size_t align) {
            void* p = std::aligned_alloc(align, /* ... */);
            if (!p) throw std::bad_alloc{};
            return p;
        }
        void deallocate(void* p, size_t bytes, size_t) {
            std::free(p);
        }
    };
  \end{minted}
  Similarly, \code{free\_list\_resource} manages a pre-allocated buffer with a free list,
  but exposes the same \code{allocate}/\code{deallocate} interface.
\end{frame}

\begin{frame}[fragile]{Polymorphic Allocator}
  The allocator holds a \code{memory\_resource*} and delegates to it.
  Since all wrapped resources share the same base, a single allocator type works with any resource.
  \begin{minted}{cpp}
    /* 2.05 */
    template <typename T>
    class polymorphic_allocator {
    public:
        using value_type = T;
        polymorphic_allocator(memory_resource* r) noexcept : resource_(r) {}

        T* allocate(size_t n) {
            return static_cast<T*>(
                resource_->allocate(n * sizeof(T), alignof(T)));
        }
        void deallocate(T* p, size_t n) noexcept {
            resource_->deallocate(p, n * sizeof(T), alignof(T));
        }
    private:
        memory_resource* resource_;
    };
  \end{minted}
\end{frame}

\begin{frame}[fragile]{Polymorphic Allocator in Action}
  \begin{minted}{cpp}
    /* 2.05 */
    int main() {
        memory_resource_wrapper<malloc_resource> malloc_res;
        alignas(16) char buf[4096];
        memory_resource_wrapper<free_list_resource> fl_res{buf, sizeof(buf)};

        polymorphic_allocator<int> a1{&malloc_res};
        polymorphic_allocator<int> a2{&fl_res};

        // Same vector type, different underlying resources!
        std::vector<int, polymorphic_allocator<int>> v1(a1);
        std::vector<int, polymorphic_allocator<int>> v2(a2);

        for (auto* v : {&v1, &v2}) {
            v->push_back(10);
            v->push_back(20);
        }
    }
  \end{minted}
\end{frame}

